\scanpage{32}
\mspage{31}
et de plus on a $\unsure{\varnothing|\phi}\in\unsure{\mathcal F|\mathfrak F|\mathscr F}_{\unsure{M|\mathcal{M}}}$ (ou encore $\mathcal F_M\neq\varnothing$).

Enfin, on voudra aussi

\uline{At$^{*}$4}$_{\unsure{0|o}}$\quad Si $\Phi,\Psi\in\deleted{\unsure{\Phi}}\mathcal F_M$, pour que $\Phi\cup\Psi$ soit \deleted{élément} dans $\mathcal F_M$ \inserted{(i.e., par At$^{*}$3, $\Phi,\Psi$ \unsure{majorés} dans $\mathcal F_M$ \inserted{pour l'inclusion})}, il f. et s. que pour tout $A\in\Phi$, $\unsure{\cdot|.}B\in\Psi$, $\Phi_A\cup\Phi_B$ soit dans $\mathcal F_M$.

Je voudrais revenir sur ce $\Phi_A$. En fait, $\Phi_A$ consiste en l'ensemble des sous-strates de $A$ dans $\Phi$, donc \uline{At$^{*}$2}$_0$ peut se formuler sous la forme plus compréhensible

\uline{At$^{*}$}2$_0$\quad Soient $\Phi$ et $\Psi$ d\deleted{\unsure{ans}}\unsure{eux} $M$-figures $\in\mathcal F_M$, et soit $A\in\Phi\cap\Psi$ une strate commune. Alors la \unsure{multistr.te|multistrate} associée \unsure{à|:} $A$ dans $\Phi$, et dans $\Psi$, est la \deleted{\unsure{m}}\unsure{même}, et elle est $\in\unsure{,|.}\mathcal F_M$.

\uline{Corollaire}\quad Si $\Phi,\Psi\in\mathcal F_M$ avec $\Phi\subset\Psi$, alors $\Phi$ est \unsure{fermée|fermé} dans $\Psi$, donc est une sous figure $(\Phi\leq\Psi)$.
